\documentclass{article}

% Keep the official style file unchanged and next to this source.
% Anonymous review mode is the default: do not add `final` or `preprint`.
\usepackage{neurips_2026}

\usepackage[T1]{fontenc}
\usepackage[utf8]{inputenc}
\usepackage{microtype}
\usepackage{booktabs}
\usepackage{amsmath}
\usepackage{xcolor}
\usepackage{url}
\usepackage{hyperref}

\definecolor{missingred}{RGB}{150,20,20}
\newcommand{\missing}[1]{\textcolor{missingred}{\textbf{[MISSING: #1]}}}
\newcommand{\system}{ClimateKG}

\title{ClimateKG: Context-Gated Scientific GraphRAG for Auditable
Land--Atmosphere Evidence Synthesis}

\author{Anonymous Author(s)}

\begin{document}
\maketitle

\begin{abstract}
Land-use and land-cover changes alter exchanges of energy, moisture, and
momentum with the atmosphere, but reported responses depend on hydroclimate,
season, atmospheric regime, terrain, and intervention geometry. Consequently,
topically relevant literature is not necessarily applicable to a local
question, and an unconstrained graph search can join individually plausible
claims supported by incompatible environments. We introduce \system{}, a
scientific GraphRAG pipeline that represents study Contexts and their Facets,
evidence-backed causal or associative Claims, land-surface Transitions,
canonical States, and exact SourceBlock provenance. At query time, it keeps
environmental similarity, coverage, intervention relevance, and mechanism
relevance distinct; gates Claims by their supporting Contexts before mechanism
search; and produces cited answers traceable to source passages. In a pilot
indexing study, the pipeline fully processed 10 heterogeneous papers into 2,184
SourceBlocks and 213 Claims, while retaining five malformed extractions for
review rather than silently repairing scientific content. \missing{Insert the
frozen query benchmark and comparative retrieval results.} The intended role
is evidence triage and hypothesis formation before expensive regional climate
experiments, not autonomous land-management prescription.
\end{abstract}

\section{Introduction}

Land surfaces influence near-surface temperature, boundary-layer development,
clouds, precipitation, and circulation by modifying albedo, roughness,
evapotranspiration, and soil moisture \citep{pielke2007,cao2020,spracklen2018}.
Yet these effects are conditional. Soil-moisture heterogeneity can organize
convection at mesoscale \citep{taylor2011}, while background wind changes where
surface-triggered convection develops \citep{froidevaux2014}. Moisture recycling
also links land conditions to remote precipitation sinks \citep{keys2012}.
Retrieval for local land--atmosphere questions must therefore ask not only
whether a passage mentions an intervention and outcome, but under which
environmental and spatial conditions its finding applies.

Passage retrieval and retrieval-augmented generation (RAG) provide useful
external evidence to language models \citep{lewis2020}. Generic graph retrieval
can expose multi-hop structure \citep{edge2024}, but topical or semantic
connectivity alone does not prevent a path from combining a dry-season result,
a wet-season mechanism, and a response observed under a different flow regime.
This is especially risky where effects are conditional, scale dependent,
contradictory, or sign reversing.

We make three contributions. First, \system{} implements a provenance-preserving
scientific representation in which every extracted object is linked to the
smallest sufficient supporting SourceBlocks. Second, it compares query and
study environments at the Facet level and gates Claims before graph traversal;
missing context remains distinct from mismatch. Third, it emits a structured
query report that exposes ranking components, rejection reasons, paths, and
citations. The central hypothesis is that context gating reduces environmentally
incompatible evidence and paths without losing relevant supported findings.
\missing{State the measured headline result after benchmark freeze.}

\begin{figure}[t]
  \centering
  \fbox{\begin{minipage}{0.95\linewidth}\centering\small
  PDFs $\rightarrow$ Nemotron Parse v1.2 $\rightarrow$ SourceBlocks
  $\rightarrow$ Contexts, Facets, Transitions, Claims
  $\rightarrow$ canonical States $\rightarrow$ embeddings $\rightarrow$ Neo4j
  \\[2pt]
  Question $\rightarrow$ Query Context $\rightarrow$ Context/Facet comparison
  $\rightarrow$ Claim gate $\rightarrow$ mechanism paths
  $\rightarrow$ SourceBlocks $\rightarrow$ grounded answer
  \end{minipage}}
  \caption{The indexing and query flows. Context gating precedes mechanism-path
  construction; all substantive answer statements retain source provenance.}
  \label{fig:pipeline}
\end{figure}

\section{Method}

\paragraph{Evidence-preserving indexing.}
Only NVIDIA Nemotron Parse v1.2 is used to convert PDFs into structured page
elements. Deterministic cleaning removes repeated page furniture and segments
the text into stable SourceBlocks while preserving sections, captions, tables,
and scientific notation. Qwen3.6:27b, served locally through Ollama, performs
only the semantic stages specified by the extraction procedure. Structured
outputs are schema validated before use. Current-paper findings are separated
from cited background, and evidence identifiers are restricted to the blocks
shown to the extraction call.

A \emph{Context} denotes a scientifically distinct environmental,
observational, or experimental regime. A child Context inherits shared Facets
from its parents. Each \emph{Facet} has a controlled scientific domain and open
notion and description; reported and spatially derived Facets retain different
provenance. A \emph{Transition} links baseline and changed Contexts. A
\emph{Claim} connects directional or categorical States, distinguishes causal
from associative relations, and is scoped to a Context or Transition. Its
\texttt{conditioning\_facet\_ids} contain only Facets that materially affect
interpretation or transfer. State concepts are normalized deterministically
with a curated alias registry; embeddings propose uncertain candidates but do
not force merges. Qwen3-Embedding-4B produces normalized 2,048-dimensional
vectors for retrieval.

\paragraph{Context comparison.}
The query parser produces optional source and target States, an intervention,
and a temporary Query Context. It does not infer environmental properties from
a place name. User-supplied or reproducibly derived Facets may enrich the Query
Context. For query Facet $q$ and candidate Facet $f$ in the same domain,
\begin{equation}
s(q,f)=\alpha\cos(E_n(q),E_n(f))+(1-\alpha)\cos(E_c(q),E_c(f)),
\end{equation}
where $E_n$ embeds domain and notion and $E_c$ also includes the description.
Directional MaxSim averages each query Facet's best same-domain match. A
weighted semantic score $S_{\mathrm{semantic}}$ uses only comparable domains,
while coverage is the fraction of query-domain weight for which the candidate
has information. Ranking uses
\begin{equation}
S_{\mathrm{context}}=S_{\mathrm{semantic}}
\left[1-\lambda_{\mathrm{missing}}(1-\mathrm{coverage})\right].
\end{equation}
Thus absence is reported and conservatively penalized but is not treated as a
contradiction.

\paragraph{Claim gating and path search.}
For a Context-scoped Claim, applicability begins with its Context score; for a
Transition-scoped Claim, it uses the baseline (FROM) Context. When conditioning
Facets are comparable, their mean best-match score is combined with Context
applicability. Unknown conditioning domains remain unknown. Claim ranking then
renormalizes configured weights over non-null applicability, semantic,
intervention, and endpoint components. Context gating is applied before beam
search. For a path of length $L$,
\begin{align}
A_{\mathrm{path}}&=\min_e A_e, &
C_{\mathrm{path}}&=\min_j C_j,\\
R_{\mathrm{path}}&=A_{\mathrm{factor}}C_{\mathrm{factor}}
\,\overline{R_e}\big/\left[1+\gamma(L-1)\right],
\end{align}
where $C_j$ is adjacent-Context coherence and unknown applicability or
coherence is flagged rather than replaced by a scientific confidence value.
Local advection, patch/edge effects, distance decay, windward/leeward effects,
moisture recycling, and remote teleconnections use separate spatial rules.
Cardinal recommendations require both local environmental/geometric data and a
transferable relation; prose alone cannot create a direction.

\paragraph{Grounded synthesis.}
After path ranking, the evidence package contains only selected Claims and their
supporting SourceBlocks. Budget trimming cannot retain a Claim after dropping
all of its evidence. The final output includes a natural-language answer and a
machine-readable report linking each substantive statement to its Claims,
supporting Contexts, Paper, and SourceBlocks. Retrieval scores rank evidence;
they are not probabilities that a scientific result is true.

\section{Pilot Validation and Evaluation Design}

\paragraph{Pilot corpus.}
We selected 15 PDFs spanning local land-use influences on weather and climate,
including irrigation, vegetation and soil-moisture heterogeneity, surface
roughness, albedo, and land-cover change. Table~\ref{tab:indexing} reports the
saved pipeline artifacts. Ten papers completed all indexing stages with no
unresolved consolidation conflicts. Five were retained as \textsc{Needs Review}
after referential or semantic-preservation checks failed. This fail-closed
behavior is intentional: malformed scientific objects are not silently passed
to retrieval. Seventeen deterministic tests cover normalization, parsing cache
identity, Context hierarchy, scope and provenance references, reconciliation,
Neo4j property serialization, missing-versus-mismatch scoring, and evidence
budget behavior.

\begin{table}[t]
\centering
\small
\caption{Saved pilot indexing artifacts. These counts establish execution and
schema validity, not retrieval superiority.}
\label{tab:indexing}
\begin{tabular}{lr}
\toprule
Artifact & Count \\
\midrule
Fully indexed / retained for review & 10 / 5 \\
SourceBlocks & 2,184 \\
Contexts / Facets & 101 / 437 \\
Transitions / Claims & 51 / 213 \\
Canonical States & 322 \\
Deterministic tests passed & 17 \\
\bottomrule
\end{tabular}
\end{table}

\paragraph{Acceptance-critical retrieval evaluation.}
\missing{Annotate 18--24 queries covering forward, backward, source-to-target,
global, contradictory, missing-context, sign-reversal, and spatial cases. State
the train/development/test split and annotator expertise.} Each gold record
should identify applicable Contexts and Claims, incompatible evidence,
acceptable mechanism paths, smallest sufficient SourceBlocks, and supported
answer propositions. A second annotator should independently review a fixed
subset, with agreement and adjudication reported.

All variants must use the same PDFs, queries, embedding model, and evidence
budget: (i) BM25 passage retrieval, (ii) dense passage retrieval, (iii) graph
retrieval without Context gating, and (iv) full \system{}. The primary metrics
are applicable-Claim precision/recall/F1, incompatible-Claim rate,
valid-path rate, citation precision and coverage, supported-answer rate, and
answer completeness. Report latency, token use, and peak memory as operational
metrics. Whole-Context-only and no-conditioning-Facet ablations isolate the
Facet-level contribution.

\begin{table}[t]
\centering
\scriptsize
\caption{Frozen benchmark results. Replace every placeholder from generated
evaluation JSON; higher is better except incompatible-Claim rate.}
\label{tab:results}
\begin{tabular}{lccccc}
\toprule
System & Claim F1 & Incompat. $\downarrow$ & Valid path & Cite. prec. & Supported \\
\midrule
BM25 & \missing{--} & \missing{--} & n/a & \missing{--} & \missing{--} \\
Dense RAG & \missing{--} & \missing{--} & n/a & \missing{--} & \missing{--} \\
Ungated graph & \missing{--} & \missing{--} & \missing{--} & \missing{--} & \missing{--} \\
\system{} & \missing{--} & \missing{--} & \missing{--} & \missing{--} & \missing{--} \\
\bottomrule
\end{tabular}
\end{table}

\missing{Add one compact qualitative example showing the same query under the
ungated and gated systems, including rejected Context, mechanism path, and exact
SourceBlock citations. Add confidence intervals or bootstrap intervals where
the benchmark size permits.}

\section{Pathway to Climate Impact and Limitations}

The intended users are climate researchers and regional-planning teams deciding
which land-surface hypotheses merit detailed investigation. Today, relevant
evidence is distributed across observational, remote-sensing, and modelling
studies whose environmental assumptions differ. \system{} can reduce the manual
search space while exposing why a finding may or may not transfer. Its output
can define a small set of evidence-backed land-cover perturbations, mechanism
hypotheses, seasons, and spatial constraints for subsequent regional climate or
weather-model experiments. This pathway may reduce expensive, weakly motivated
scenario runs and make the literature basis of each tested scenario auditable.
It does not establish that an intervention will produce a desired local outcome.

Potential harms arise if retrieval rank is mistaken for causal certainty or if
literature coverage reproduces geographic and institutional biases. A locally
plausible climate response may still conflict with biodiversity, water, food,
livelihood, or justice objectives. Accordingly, the system must expose missing
context, contradictory findings, provenance, and uncertainty; local experts and
affected communities remain responsible for framing objectives and evaluating
trade-offs. \missing{Add feedback from at least one land--atmosphere domain
expert and, if recommendations are discussed, one intended-user or stakeholder
review.}

The present pilot is small and not a representative census of land--atmosphere
science. Extraction quality depends on PDF parsing and local language-model
judgment. Five fail-closed papers show that structured extraction remains a
material bottleneck. Spatial enrichment is incomplete, and no WRF or field
validation has been performed. \missing{After completing the benchmark, add
measured failure categories, generalization limits, and compute/energy use.}

\section{Conclusion}

\system{} makes environmental applicability an explicit part of scientific
GraphRAG. It preserves exact evidence, separates missing information from
mismatch, selects Claims before path construction, and prevents spatial prose
from becoming unsupported local direction. The pilot demonstrates traceable,
validated indexing with visible failure handling. \missing{Conclude with the
measured retrieval and grounding result, without claiming scientific truth or
management effectiveness.}

\bibliographystyle{plainnat}
\bibliography{references}

\end{document}
